% Generated draft: requires mathematical and source review.
\documentclass[11pt]{article}
\usepackage[margin=25mm]{geometry}
\usepackage{fontspec}
\setmainfont{Latin Modern Roman}
\usepackage{amsmath,amssymb,mathtools,unicode-math}
\setmathfont{Latin Modern Math}
\usepackage{enumitem,needspace,xurl,hyperref}
\setlength{\parindent}{0pt}
\setlength{\parskip}{5pt}
\setlength{\emergencystretch}{3em}
\newcommand{\N}{\mathbb{N}}
\newcommand{\Z}{\mathbb{Z}}
\newcommand{\Q}{\mathbb{Q}}
\newcommand{\R}{\mathbb{R}}
\newcommand{\C}{\mathbb{C}}
\newcommand{\F}{\mathbb{F}}
\newcommand{\A}{\mathbb{A}}
\newcommand{\PP}{\mathbb{P}}
\newcommand{\E}{\mathbb{E}}
\newcommand{\eps}{\varepsilon}
\newcommand{\dd}{\,\mathrm{d}}
\begin{document}
\section*{UnsolvedMath ID: 700402}
\textbf{Source problem code:} catalogue entry
\subsection{Definitions and mathematical statement}
Let \(G\) be a group with identity \(1\). The group is finitely generated if \(G=\langle S\rangle\) for some finite subset \(S\subseteq G\), simple if its only normal subgroups are \(\{1\}\) and \(G\) and \(G\ne\{1\}\), and periodic if every element has finite order. For an element of finite order, write
\[
\operatorname{ord}(g)=\min\{n\in\Z_{\geq 1}:g^n=1\}.
\]
An involution is an element of order \(2\); a nontrivial element of odd order has order \(m\) for some odd integer \(m\geq3\).

The supplied statement uses the expression ``(conjugacy) biprimitively finite'' without defining it. Let \(\mathcal P(G)\) stand provisionally for the intended property. The two apparent readings are ``biprimitively finite'' and ``conjugacy biprimitively finite''; the record does not specify their definitions or whether both existence questions are intended.

Subject to resolving that qualification, the question attributed to V. P. Shunkov is whether there exists a group \(G\) such that
\[
\begin{gathered}
|G|=\infty,\qquad
\exists S\subseteq G\;\bigl(|S|<\infty\ \land\ G=\langle S\rangle\bigr),\\
G\text{ is simple},\qquad
\forall g\in G\ \exists n\in\Z_{\geq1}\quad g^n=1,\qquad
\mathcal P(G),\\
\exists x,y\in G\ \exists m\in\Z_{\geq3}\quad
\bigl(m\text{ is odd}\ \land\ \operatorname{ord}(x)=2\ \land\ \operatorname{ord}(y)=m\bigr).
\end{gathered}
\]
Here \(\mathcal P\) is an unresolved placeholder, not a definition of either finiteness condition. This transcription therefore specifies the remaining hypotheses and quantifiers but requires clarification of the named property before it is a complete mathematical formulation.
\subsection{Short English statement}
Does there exist an infinite, finitely generated, simple periodic group satisfying the source's ``(conjugacy) biprimitively finite'' condition and containing both an involution and a nonidentity element of odd order? The named finiteness condition and the intended role of the parenthesis require clarification.
\subsection{Sources}
\begin{itemize}
\item Local UnsolvedMath registry, supplied complete source record: permanent ID \(700402\), exact source code \texttt{catalogue entry}, title ``Kourovka Notebook Problem 6.60''. Supplied upstream revision: \texttt{d371ee917ac9647fbb06cd6fc529deec6f3dda82}. The registry labels the problem \texttt{open}; this is supplied source metadata, not an independently verified assessment.
\item The supplied attribution is to V. P. Shunkov, proposed in 1978. The cited source is \emph{The Kourovka Notebook (Unsolved Problems in Group Theory)}, original problem number 6.60, 6th issue (1978), as represented in the 21st issue edition (2026), \texttt{arXiv:1401.0300v48}. The supplied collection description names E. I. Khukhro and V. D. Mazurov as editors of the Issues 1--21 collection (1965--2026). These bibliographic details are reproduced from the record without independent source verification.
\end{itemize}
\end{document}
